% Blueprint: Quaternion Algebras
% Created by Numina

\begin{document}

\section{Quaternion algebras: beginnings}

Throughout, $F$ is a commutative field with $\operatorname{char} F \neq 2$, and
$a, b \in F^\times$. We follow Voight, \emph{Quaternion Algebras}, Chapter~2.

In Mathlib the quaternion algebra is $\mathbb{H}[R, c_1, c_2, c_3]$ with relations
$i^2 = c_1 + c_2 i$, $j^2 = c_3$, $ij = k$ and $ji = c_2 j - k$. Voight's
$\left(\frac{a,b}{F}\right)$ is the special case $c_2 = 0$, i.e.
$\mathbb{H}[F, a, 0, b]$, for which $i^2 = a$, $j^2 = b$, $ij = k$, $ji = -k = -ij$.

\subsection{The algebra and its multiplication}

\begin{definition}[Quaternion algebra]
    \label{def:quaternion-algebra}
    For $a, b \in F^\times$, the \emph{quaternion algebra}
    $\left(\frac{a,b}{F}\right)$ is the $F$-algebra $\mathbb{H}[F, a, 0, b]$ with
    underlying $F$-vector space spanned by $1, i, j, k$ subject to
    $i^2 = a$, $j^2 = b$, $ji = -ij$, and $k := ij$.
\end{definition}

\begin{definition}[Standard generators]
    \label{def:standard-generators}
    \uses{def:quaternion-algebra}
    The \emph{standard generators} of $\left(\frac{a,b}{F}\right)$ are the elements
    $i = (0,1,0,0)$ and $j = (0,0,1,0)$, together with the derived element
    $k = ij = (0,0,0,1)$. They satisfy $i^2 = a$, $j^2 = b$, $ij = k$, $ji = -k$.
\end{definition}

\begin{lemma}[Anticommutation of generators]
    \label{lem:ji-eq-neg-ij}
    \uses{def:standard-generators}
    In $\left(\frac{a,b}{F}\right)$ we have $ji = -ij$.
\end{lemma}
\begin{proof}
    \uses{def:standard-generators}
    Immediate from the defining relation $ji = c_2 j - k$ with $c_2 = 0$ and
    $ij = k$, i.e. the Mathlib basis relation \texttt{j\_mul\_i} specialized to
    $c_2 = 0$.
\end{proof}

\begin{lemma}[Square of $k$]
    \label{lem:k-squared}
    \uses{def:standard-generators, lem:ji-eq-neg-ij}
    With $k = ij$ we have $k^2 = -ab$.
\end{lemma}
\begin{proof}
    \uses{def:standard-generators, lem:ji-eq-neg-ij}
    Compute $(ij)(ij) = i(ji)j = i(-ij)j = -(i^2)(j^2) = -ab$, using
    associativity, $ji = -ij$ (\ref{lem:ji-eq-neg-ij}), and $i^2 = a$,
    $j^2 = b$. Each rewrite is an instance of the Mathlib multiplication lemmas
    for $\mathbb{H}[F,a,0,b]$ (\texttt{imK\_mul}, \texttt{imI\_mul}).
\end{proof}

\begin{lemma}[Mixed products]
    \label{lem:mul-table}
    \uses{def:standard-generators, lem:ji-eq-neg-ij}
    With $k = ij$ we have $ik = aj$, $ki = -aj$, $jk = -bi$, and $kj = bi$.
\end{lemma}
\begin{proof}
    \uses{def:standard-generators, lem:ji-eq-neg-ij}
    For example $j k = j(ij) = (ji)j = (-ij)j = -i(j^2) = -bi$, using
    \ref{lem:ji-eq-neg-ij} and $j^2 = b$; the others are analogous and reduce to
    the Mathlib basis product lemmas \texttt{i\_mul\_k}, \texttt{k\_mul\_i},
    \texttt{j\_mul\_k}, \texttt{k\_mul\_j} with $c_2 = 0$.
\end{proof}

\begin{theorem}[Dimension four]
    \label{thm:dim-four}
    \uses{def:quaternion-algebra}
    $\left(\frac{a,b}{F}\right)$ has dimension $4$ as an $F$-vector space.
\end{theorem}
\begin{proof}
    \uses{def:quaternion-algebra}
    This is Mathlib's \texttt{QuaternionAlgebra.finrank\_eq\_four}: the linear
    equivalence with $F^4$ (sending a quaternion to its four coordinates) shows the
    algebra is free of rank $4$.
\end{proof}

\begin{lemma}[Standard basis]
    \label{lem:standard-basis}
    \uses{def:standard-generators, thm:dim-four}
    The elements $1, i, j, k$ form an $F$-basis of $\left(\frac{a,b}{F}\right)$.
    In particular they are $F$-linearly independent: if
    $t + xi + yj + zk = 0$ with $t,x,y,z \in F$, then $t = x = y = z = 0$.
\end{lemma}
\begin{proof}
    \uses{thm:dim-four}
    Under the coordinate linear equivalence
    $\mathbb{H}[F,a,0,b] \simeq F^4$ (Mathlib \texttt{linearEquivTuple}) the
    elements $1,i,j,k$ map to the standard basis $e_0, e_1, e_2, e_3$ of $F^4$,
    which is linearly independent and spanning; transport along the equivalence.
\end{proof}

\begin{lemma}[Lifting generators to an algebra homomorphism]
    \label{lem:lifthom-construction}
    \uses{def:standard-generators}
    Let $B$ be an $F$-algebra containing elements $i', j'$ with $(i')^2 = a$,
    $(j')^2 = b$, and $j'i' = -i'j'$. Then $(i', j', i'j')$ form a
    \texttt{QuaternionAlgebra.Basis} for $B$ with parameters $(a, 0, b)$, and
    \texttt{QuaternionAlgebra.Basis.liftHom} produces an $F$-algebra homomorphism
    $\varphi \colon \left(\frac{a,b}{F}\right) \to B$ with $\varphi(i) = i'$,
    $\varphi(j) = j'$, and $\varphi(k) = i'j'$.
\end{lemma}
\begin{proof}
    \uses{def:standard-generators}
    The three relations $(i')^2 = a = c_1$, $(j')^2 = b = c_3$, and
    $j'i' = -i'j' = c_2 j' - i'j'$ with $c_2 = 0$ are exactly the data of a
    \texttt{QuaternionAlgebra.Basis} with parameters $(a,0,b)$. Its universal
    property \texttt{QuaternionAlgebra.Basis.liftHom} returns the asserted
    $F$-algebra homomorphism, with the images of $i, j, k$ read off from
    \texttt{QuaternionAlgebra.Basis.liftHom\_apply}.
\end{proof}

\begin{lemma}[Standard generators presentation]
    \label{lem:standard-generators-presentation}
    \uses{def:standard-generators, lem:standard-basis, lem:lifthom-construction}
    If an $F$-algebra $B$ contains nonzero $i', j'$ generating $B$ with
    $(i')^2 = a$, $(j')^2 = b$, $i'j' = -j'i'$ for $a, b \in F^\times$, then
    $1, i', j', i'j'$ is an $F$-basis of $B$, so $B \cong \left(\frac{a,b}{F}\right)$.
\end{lemma}
\begin{proof}
    \uses{lem:standard-basis, lem:lifthom-construction, thm:dim-four}
    By \ref{lem:lifthom-construction} the relations give an $F$-algebra
    homomorphism $\varphi \colon \left(\frac{a,b}{F}\right) \to B$, surjective
    because $i', j'$ generate $B$ (\texttt{QuaternionAlgebra.Basis.range\_liftHom}).
    Its source and target both have dimension $4$ (\ref{thm:dim-four},
    \ref{lem:standard-basis}), so by
    \texttt{LinearMap.injective\_iff\_surjective\_of\_finrank\_eq\_finrank} the
    surjection $\varphi$ is also injective, hence an isomorphism. Linear
    independence of $1, i', j', i'j'$ follows by transporting \ref{lem:standard-basis}.
\end{proof}

\begin{lemma}[Commuting with $i$ kills the $j,k$ parts]
    \label{lem:central-commute-i}
    \uses{lem:standard-basis, lem:mul-table, lem:ji-eq-neg-ij}
    If $\alpha = t + xi + yj + zk$ satisfies $\alpha i = i\alpha$, then $y = z = 0$.
\end{lemma}
\begin{proof}
    \uses{lem:standard-basis, lem:mul-table, lem:ji-eq-neg-ij}
    Using $i^2 = a$, $ji = -k$ and $ki = -aj$ (\ref{lem:ji-eq-neg-ij},
    \ref{lem:mul-table}) one gets $\alpha i = xa + ti - zaj - yk$, while
    $ij = k$, $ik = aj$ give $i\alpha = xa + ti + zaj + yk$. Hence
    $\alpha i - i\alpha = -2za\,j - 2y\,k = 0$. Comparing the $j$- and
    $k$-coordinates via \ref{lem:standard-basis} and using $a \neq 0$,
    $\operatorname{char} F \neq 2$ yields $z = 0$ and $y = 0$.
\end{proof}

\begin{lemma}[Commuting with $j$ kills the $i$ part]
    \label{lem:central-commute-j}
    \uses{lem:standard-basis, lem:mul-table, lem:ji-eq-neg-ij}
    If $\alpha = t + xi + yj + zk$ satisfies $\alpha j = j\alpha$, then $x = 0$.
\end{lemma}
\begin{proof}
    \uses{lem:standard-basis, lem:mul-table, lem:ji-eq-neg-ij}
    Using $j^2 = b$, $ij = k$ and $kj = bi$ (\ref{lem:ji-eq-neg-ij},
    \ref{lem:mul-table}) one gets $\alpha j = yb + zb\,i + tj + x\,k$, while
    $ji = -k$, $jk = -bi$ give $j\alpha = yb - zb\,i + tj - x\,k$. Hence
    $\alpha j - j\alpha = 2zb\,i + 2x\,k = 0$, and comparing the $i$-coordinate
    via \ref{lem:standard-basis} with $b \neq 0$, $\operatorname{char} F \neq 2$
    gives $x = 0$.
\end{proof}

\begin{theorem}[Centrality]
    \label{thm:central}
    \uses{def:quaternion-algebra, lem:standard-basis, lem:central-commute-i, lem:central-commute-j, lem:ji-eq-neg-ij}
    $\left(\frac{a,b}{F}\right)$ is central over $F$: its center equals the image
    of $F$.
\end{theorem}
\begin{proof}
    \uses{lem:standard-basis, lem:central-commute-i, lem:central-commute-j}
    By \texttt{Algebra.IsCentral.mem\_center\_iff} it suffices to show every
    central $\alpha$ lies in the image of $F$. Write
    $\alpha = t + xi + yj + zk$ (\ref{lem:standard-basis}). Centrality gives
    $\alpha i = i\alpha$, so $y = z = 0$ (\ref{lem:central-commute-i}), and
    $\alpha j = j\alpha$, so $x = 0$ (\ref{lem:central-commute-j}). Hence
    $\alpha = t \in F$.
\end{proof}

\subsection{Matrix representations}

\begin{definition}[Matrix basis over a splitting field]
    \label{def:matrix-basis}
    \uses{def:quaternion-algebra}
    Let $K = F(\sqrt a)$ contain a root $s$ of $x^2 - a$. Define elements of
    $M_2(K)$ by
    \[
      I = \begin{pmatrix} s & 0 \\ 0 & -s \end{pmatrix}, \qquad
      J = \begin{pmatrix} 0 & b \\ 1 & 0 \end{pmatrix}, \qquad K_0 = IJ .
    \]
    These satisfy $I^2 = a\cdot 1$, $J^2 = b \cdot 1$, $IJ = K_0$,
    $JI = -K_0$, forming a \texttt{QuaternionAlgebra.Basis} with parameters
    $(a, 0, b)$.
\end{definition}

\begin{lemma}[Matrix basis relations]
    \label{lem:matrix-basis-relations}
    \uses{def:matrix-basis}
    The matrices $I, J$ of \ref{def:matrix-basis} satisfy
    $I^2 = a\cdot 1$, $J^2 = b\cdot 1$, and $JI = -IJ$ in $M_2(K)$.
\end{lemma}
\begin{proof}
    \uses{def:matrix-basis}
    Direct $2\times 2$ matrix multiplication using $s^2 = a$; each entry equality
    is checked by \texttt{Matrix} simp lemmas (\texttt{Matrix.mul\_fin\_two}).
\end{proof}

\begin{lemma}[Matrix liftHom]
    \label{lem:matrix-lifthom}
    \uses{def:matrix-basis, lem:matrix-basis-relations, lem:lifthom-construction}
    There is an $F$-algebra homomorphism
    $\lambda \colon \left(\frac{a,b}{F}\right) \to M_2(F(\sqrt a))$ with
    $\lambda(i) = I$, $\lambda(j) = J$, $\lambda(k) = IJ$.
\end{lemma}
\begin{proof}
    \uses{lem:matrix-basis-relations, lem:lifthom-construction}
    By \ref{lem:matrix-basis-relations} the matrices $I, J$ satisfy
    $I^2 = a$, $J^2 = b$, $JI = -IJ$ in $M_2(K)$, which are the hypotheses of
    \ref{lem:lifthom-construction} with $B = M_2(F(\sqrt a))$. Applying it gives
    the $F$-algebra homomorphism $\lambda$ with the stated images.
\end{proof}

\begin{proposition}[Matrix embedding]
    \label{prop:matrix-embedding}
    \uses{def:matrix-basis, lem:matrix-lifthom, thm:dim-four, lem:standard-basis}
    The $F$-algebra homomorphism
    $\lambda \colon \left(\frac{a,b}{F}\right) \to M_2(F(\sqrt a))$ of
    \ref{lem:matrix-lifthom} is injective and an isomorphism onto its image.
\end{proposition}
\begin{proof}
    \uses{lem:matrix-lifthom, lem:standard-basis}
    By \ref{lem:matrix-lifthom} we have the $F$-algebra homomorphism $\lambda$.
    For injectivity: if $\lambda(t + xi + yj + zk) = 0$ then the matrix
    $\left(\begin{smallmatrix} t + xs & b(y+zs) \\ y - zs & t - xs\end{smallmatrix}\right) = 0$;
    reading off entries and using linear independence of $1, s$ over $F$ (as $K/F$
    is a field extension, or $s \in F$) gives $t = x = y = z = 0$, so by
    \ref{lem:standard-basis} the kernel is trivial.
\end{proof}

\begin{corollary}[Split quaternion algebra is a matrix algebra]
    \label{cor:split-matrix}
    \uses{prop:matrix-embedding, thm:dim-four}
    There is an $F$-algebra isomorphism
    $\left(\frac{1,b}{F}\right) \xrightarrow{\sim} M_2(F)$, given by
    $i \mapsto \left(\begin{smallmatrix}1 & 0\\ 0 & -1\end{smallmatrix}\right)$,
    $j \mapsto \left(\begin{smallmatrix}0 & b\\ 1 & 0\end{smallmatrix}\right)$.
    In particular $\left(\frac{1,1}{F}\right) \cong M_2(F)$.
\end{corollary}
\begin{proof}
    \uses{prop:matrix-embedding, thm:dim-four}
    Specialize \ref{prop:matrix-embedding} with $a = 1$, $\sqrt a = 1 \in F$, so
    $K = F$ and $\lambda$ lands in $M_2(F)$. It is an injective $F$-algebra
    homomorphism between spaces of equal dimension
    $\dim_F \left(\frac{1,b}{F}\right) = 4 = \dim_F M_2(F)$
    (\ref{thm:dim-four}), hence surjective by
    \texttt{LinearMap.injective\_iff\_surjective\_of\_finrank\_eq\_finrank}, hence
    bijective; promote to an $F$-algebra isomorphism via
    \texttt{AlgEquiv.ofBijective}.
\end{proof}

\subsection{Conjugation, norm, and pure quaternions}

\begin{definition}[Standard involution]
    \label{def:conjugation}
    \uses{def:quaternion-algebra}
    The \emph{conjugation} (standard involution) on $\left(\frac{a,b}{F}\right)$ is
    the map $\alpha = t + xi + yj + zk \mapsto \bar\alpha = t - xi - yj - zk$,
    realized in Mathlib by the \texttt{star} operation.
\end{definition}

\begin{lemma}[Trace and norm via conjugation]
    \label{lem:trace-norm}
    \uses{def:conjugation, lem:standard-basis}
    For $\alpha = t + xi + yj + zk$ in the Hamiltonians
    $\mathbb{H} = \left(\frac{-1,-1}{\mathbb{R}}\right)$,
    \[
      \alpha + \bar\alpha = 2t, \qquad
      \alpha \bar\alpha = \bar\alpha \alpha = t^2 + x^2 + y^2 + z^2 \in \mathbb{R}.
    \]
\end{lemma}
\begin{proof}
    \uses{def:conjugation}
    The first identity is \texttt{Quaternion.add\_star}. The second is
    \texttt{Quaternion.self\_mul\_star} / \texttt{star\_mul\_self}, which equal the
    coercion of $\operatorname{normSq}\alpha = t^2 + x^2 + y^2 + z^2$
    (\texttt{normSq\_def}).
\end{proof}

\begin{lemma}[Multiplicativity of the norm]
    \label{lem:norm-mul}
    \uses{lem:trace-norm}
    For the squared norm $\lVert\alpha\rVert^2 = \operatorname{normSq}\alpha = \alpha\bar\alpha$
    on $\mathbb{H}$, we have
    $\lVert\alpha\beta\rVert^2 = \lVert\alpha\rVert^2 \lVert\beta\rVert^2$.
\end{lemma}
\begin{proof}
    \uses{lem:trace-norm}
    $\operatorname{normSq}$ is the monoid-with-zero homomorphism
    $\texttt{Quaternion.normSq} \colon \mathbb{H} \to^{*}_{0} \mathbb{R}$, so
    multiplicativity is \texttt{map\_mul}.
\end{proof}

\begin{definition}[Pure (imaginary) quaternion]
    \label{def:pure}
    \uses{def:conjugation}
    A quaternion $\alpha$ is \emph{pure} (imaginary) if its real part vanishes,
    equivalently $\bar\alpha = -\alpha$. The pure quaternions form the subspace
    $\mathbb{H}^0 = \mathbb{R}i + \mathbb{R}j + \mathbb{R}k \cong \mathbb{R}^3$.
\end{definition}

\begin{lemma}[Square of a pure quaternion]
    \label{lem:pure-square}
    \uses{def:pure}
    A quaternion $v$ is pure if and only if $v^2 = -\lVert v\rVert^2$. In particular
    for pure $v$, $v^2 = -\lVert v\rVert^2 \in \mathbb{R}_{\le 0}$.
\end{lemma}
\begin{proof}
    \uses{def:pure}
    This is Mathlib's \texttt{Quaternion.sq\_eq\_neg\_normSq}:
    $v^2 = -\operatorname{normSq} v \iff v.\mathrm{re} = 0$.
\end{proof}

\begin{lemma}[Real part of a product of pure quaternions]
    \label{lem:pure-product-re}
    \uses{def:pure, lem:standard-basis, lem:ji-eq-neg-ij, lem:k-squared, lem:mul-table}
    For pure quaternions $v = v_1 i + v_2 j + v_3 k$ and
    $w = w_1 i + w_2 j + w_3 k$ in $\mathbb{H}^0$,
    \[
      \operatorname{re}(vw) = -(v_1 w_1 + v_2 w_2 + v_3 w_3) = -(v \cdot w).
    \]
\end{lemma}
\begin{proof}
    \uses{lem:standard-basis, lem:ji-eq-neg-ij, lem:k-squared, lem:mul-table}
    Expand $vw$ by bilinearity over the basis $i, j, k$ and apply the
    multiplication table (\ref{lem:ji-eq-neg-ij}, \ref{lem:k-squared},
    \ref{lem:mul-table}, with $a = b = -1$): the real contributions come only from
    $i^2 = j^2 = k^2 = -1$, giving real part
    $-(v_1 w_1 + v_2 w_2 + v_3 w_3)$. The mixed terms are pure and do not
    contribute to $\operatorname{re}$ (\ref{lem:standard-basis}).
\end{proof}

\begin{lemma}[Imaginary part of a product of pure quaternions]
    \label{lem:pure-product-im}
    \uses{def:pure, lem:standard-basis, lem:ji-eq-neg-ij, lem:k-squared, lem:mul-table}
    For pure quaternions $v = v_1 i + v_2 j + v_3 k$ and
    $w = w_1 i + w_2 j + w_3 k$ in $\mathbb{H}^0$, the imaginary part of $vw$ is
    the cross product:
    \[
      \operatorname{im}(vw) = v \times w
      = (v_2 w_3 - v_3 w_2)\,i + (v_3 w_1 - v_1 w_3)\,j + (v_1 w_2 - v_2 w_1)\,k.
    \]
\end{lemma}
\begin{proof}
    \uses{lem:standard-basis, lem:ji-eq-neg-ij, lem:k-squared, lem:mul-table}
    Expanding $vw$ over the basis (\ref{lem:ji-eq-neg-ij}, \ref{lem:k-squared},
    \ref{lem:mul-table}, $a=b=-1$), the off-diagonal terms $ij = k$, $jk = i$,
    $ki = j$ and their anticommuted partners collect the $i,j,k$ coordinates to
    exactly the entries of \texttt{crossProduct} $v\,w$ (\texttt{cross\_apply});
    compare coordinates via \ref{lem:standard-basis}.
\end{proof}

\begin{lemma}[Product of pure quaternions]
    \label{lem:pure-product}
    \uses{def:pure, lem:standard-basis, lem:pure-product-re, lem:pure-product-im, lem:ji-eq-neg-ij, lem:k-squared, lem:mul-table}
    For pure quaternions $v = v_1 i + v_2 j + v_3 k$ and
    $w = w_1 i + w_2 j + w_3 k$ in $\mathbb{H}^0 \cong \mathbb{R}^3$,
    \[
      vw = -(v \cdot w) + (v \times w),
    \]
    where $v\cdot w$ is the dot product and $v \times w$ the cross product (a pure
    quaternion).
\end{lemma}
\begin{proof}
    \uses{lem:pure-product-re, lem:pure-product-im}
    Split $vw$ into its real and imaginary parts. The real part is $-(v\cdot w)$
    by \ref{lem:pure-product-re} and the imaginary part is $v\times w$ by
    \ref{lem:pure-product-im}; adding them gives the claim.
\end{proof}

\begin{lemma}[Orthogonality criteria for pure quaternions]
    \label{lem:pure-orthogonal}
    \uses{lem:pure-product}
    For pure quaternions $v, w$:
    (a) $vw$ is pure if and only if $v, w$ are orthogonal;
    (b) $wv = -vw$ if and only if $v, w$ are orthogonal.
\end{lemma}
\begin{proof}
    \uses{lem:pure-product}
    By \ref{lem:pure-product}, $vw = -(v\cdot w) + v\times w$ and
    $wv = -(v\cdot w) - v\times w$ (the cross product is antisymmetric). Thus the
    real part of $vw$ is $-(v\cdot w)$, vanishing iff $v\cdot w = 0$, giving (a); and
    $wv + vw = -2(v\cdot w)$, so $wv = -vw$ iff $v\cdot w = 0$, giving (b).
\end{proof}

\subsection{Rotations}

We work with the real Hamiltonians $\mathbb{H} = \left(\frac{-1,-1}{\mathbb{R}}\right)$,
the unit group $\mathbb{H}^1 = \{\alpha : \lVert\alpha\rVert^2 = 1\}$, and
$\mathbb{H}^0 \cong \mathbb{R}^3$ the pure quaternions.

\begin{definition}[Euclidean structure on the pure quaternions]
    \label{def:pure-euclidean}
    \uses{def:pure, lem:trace-norm}
    Let $\mathbb{H}^0 \subseteq \mathbb{H}$ be the $\mathbb{R}$-submodule of pure
    quaternions. Equip it with the inner product
    $\langle v, w\rangle = \tfrac12(\bar v w + \bar w v) \in \mathbb{R}$ inherited
    from $\operatorname{normSq}$ (\ref{lem:trace-norm}), for which $i, j, k$ is an
    orthonormal basis. The resulting map
    $e \colon \mathbb{H}^0 \xrightarrow{\sim} \mathrm{EuclideanSpace}\ \mathbb{R}\ (\mathrm{Fin}\,3)$
    sending $v_1 i + v_2 j + v_3 k \mapsto (v_1, v_2, v_3)$ is a linear isometric
    equivalence (\texttt{OrthonormalBasis}, \texttt{LinearEquiv.isometryOfOrthonormal}).
\end{definition}

\begin{lemma}[A norm-preserving map gives a special orthogonal matrix]
    \label{lem:isometry-to-so}
    \uses{def:pure-euclidean}
    Let $\rho \colon \mathbb{H}^0 \to \mathbb{H}^0$ be an $\mathbb{R}$-linear map
    with $\lVert \rho(v)\rVert = \lVert v\rVert$ for all $v$ and $\det \rho = 1$.
    Then, transported across the isometry $e$ of \ref{def:pure-euclidean}, the
    matrix of $\rho$ in the orthonormal basis $i, j, k$ lies in
    $\mathrm{SO}(3) = \texttt{Matrix.specialOrthogonalGroup}\ (\mathrm{Fin}\,3)\ \mathbb{R}$.
\end{lemma}
\begin{proof}
    \uses{def:pure-euclidean}
    A norm-preserving $\mathbb{R}$-linear endomorphism of a finite-dimensional
    inner product space is a linear isometry, hence (transported by $e$) a linear
    isometry of $\mathrm{EuclideanSpace}\ \mathbb{R}\ (\mathrm{Fin}\,3)$; its matrix
    is orthogonal by \texttt{EuclideanGroup.linearIsometryEquivToOrthogonal}, so it
    satisfies $A^\top A = 1$ (\texttt{Matrix.mem\_orthogonalGroup\_iff}). Together
    with $\det A = \det\rho = 1$, the criterion
    \texttt{Matrix.mem\_specialOrthogonalGroup\_iff} places $A$ in
    $\mathrm{SO}(3)$.
\end{proof}

\begin{lemma}[Conjugation preserves purity]
    \label{lem:conj-preserves-pure}
    \uses{def:pure, lem:trace-norm}
    For $\alpha \in \mathbb{H}^\times$ and pure $v$, the element
    $\alpha v \alpha^{-1}$ is pure.
\end{lemma}
\begin{proof}
    \uses{def:pure, lem:trace-norm}
    $v$ is pure iff $v + \bar v = 0$ (\ref{def:pure}). Since conjugation is an
    anti-automorphism with
    $\overline{\alpha v \alpha^{-1}} = \overline{\alpha^{-1}}\,\bar v\,\bar\alpha$
    and $\bar v = -v$, one computes
    $\alpha v \alpha^{-1} + \overline{\alpha v \alpha^{-1}} = 0$ using
    $\bar\alpha = \lVert\alpha\rVert^2 \alpha^{-1}$ (\ref{lem:trace-norm}). Hence the
    real part vanishes.
\end{proof}

\begin{lemma}[Conjugation preserves the norm]
    \label{lem:conj-preserves-norm}
    \uses{lem:norm-mul}
    For $\alpha \in \mathbb{H}^1$ and any $v$,
    $\lVert \alpha v \alpha^{-1}\rVert^2 = \lVert v\rVert^2$.
\end{lemma}
\begin{proof}
    \uses{lem:norm-mul}
    By multiplicativity (\ref{lem:norm-mul}),
    $\lVert\alpha v \alpha^{-1}\rVert^2
    = \lVert\alpha\rVert^2 \lVert v\rVert^2 \lVert\alpha^{-1}\rVert^2$, and
    $\lVert\alpha\rVert^2 = 1$ with
    $\lVert\alpha^{-1}\rVert^2 = \lVert\alpha\rVert^{-2} = 1$ for
    $\alpha \in \mathbb{H}^1$.
\end{proof}

\begin{lemma}[Conjugation is linear on pure quaternions]
    \label{lem:conj-linear}
    \uses{lem:conj-preserves-pure}
    For fixed $\alpha \in \mathbb{H}^1$, the map
    $\rho_\alpha \colon \mathbb{H}^0 \to \mathbb{H}^0$,
    $v \mapsto \alpha v \alpha^{-1}$, is $\mathbb{R}$-linear (well defined into
    $\mathbb{H}^0$ by \ref{lem:conj-preserves-pure}).
\end{lemma}
\begin{proof}
    \uses{lem:conj-preserves-pure}
    Left and right multiplication in the algebra $\mathbb{H}$ are
    $\mathbb{R}$-linear, so $v \mapsto \alpha v \alpha^{-1}$ is $\mathbb{R}$-linear;
    \ref{lem:conj-preserves-pure} ensures the codomain is $\mathbb{H}^0$.
\end{proof}

\begin{lemma}[$\rho_\alpha$ is orthogonal]
    \label{lem:rotation-orthogonal}
    \uses{lem:conj-preserves-norm, lem:conj-linear, def:pure-euclidean}
    For $\alpha \in \mathbb{H}^1$, the map $\rho_\alpha$ of \ref{lem:conj-linear}
    preserves the norm on $\mathbb{H}^0$: $\lVert \rho_\alpha(v)\rVert = \lVert v\rVert$
    for all pure $v$.
\end{lemma}
\begin{proof}
    \uses{lem:conj-preserves-norm}
    This is exactly \ref{lem:conj-preserves-norm} read for the squared norm
    $\lVert\alpha v\alpha^{-1}\rVert^2 = \lVert v\rVert^2$, taking square roots
    (both norms are nonnegative).
\end{proof}

\begin{lemma}[$\det \rho_\alpha = 1$]
    \label{lem:rotation-det-one}
    \uses{lem:conj-linear, def:pure-euclidean}
    For $\alpha \in \mathbb{H}^1$, the determinant of the $3\times 3$ matrix of
    $\rho_\alpha$ in the orthonormal basis $i, j, k$ equals
    $(\operatorname{normSq}\alpha)^3 = 1$.
\end{lemma}
\begin{proof}
    \uses{lem:conj-linear}
    Write $\alpha = t + xi + yj + zk$ and compute $\rho_\alpha(v) = \alpha v\bar\alpha/\operatorname{normSq}\alpha$
    on each basis vector; the resulting explicit $3\times 3$ matrix is the standard
    quaternion rotation matrix whose determinant is $(\operatorname{normSq}\alpha)^3$
    by a direct determinant computation. Since $\alpha \in \mathbb{H}^1$ means
    $\operatorname{normSq}\alpha = 1$, the determinant is $1$.
\end{proof}

\begin{proposition}[Conjugation acts by rotations]
    \label{prop:rotation}
    \uses{lem:rotation-orthogonal, lem:rotation-det-one, lem:isometry-to-so, lem:conj-linear}
    For $\alpha \in \mathbb{H}^1$, the linear map $\rho_\alpha$ of
    \ref{lem:conj-linear} on $\mathbb{H}^0 \cong \mathbb{R}^3$ lies in
    $\mathrm{SO}(3)$.
\end{proposition}
\begin{proof}
    \uses{lem:rotation-orthogonal, lem:rotation-det-one, lem:isometry-to-so}
    By \ref{lem:rotation-orthogonal} the map $\rho_\alpha$ preserves the norm and
    by \ref{lem:rotation-det-one} its determinant is $1$. The bridge
    \ref{lem:isometry-to-so} then places its matrix in
    $\mathrm{SO}(3) = \texttt{Matrix.specialOrthogonalGroup}\ (\mathrm{Fin}\,3)\ \mathbb{R}$.
\end{proof}

\begin{lemma}[Kernel of the rotation map]
    \label{lem:rotation-kernel}
    \uses{prop:rotation, lem:standard-basis, thm:central}
    If $\alpha \in \mathbb{H}^1$ satisfies $\alpha v \alpha^{-1} = v$ for all pure
    $v$, then $\alpha = \pm 1$.
\end{lemma}
\begin{proof}
    \uses{prop:rotation, lem:standard-basis, lem:pure-orthogonal, thm:central}
    The hypothesis says $\alpha$ commutes with every pure quaternion, hence with all
    of $\mathbb{H}$, so $\alpha$ is central; by centrality (the real analogue of
    \ref{thm:central}) $\alpha \in \mathbb{R}$. Then
    $\lVert\alpha\rVert^2 = \alpha^2 = 1$ gives $\alpha = \pm 1$.
\end{proof}

\begin{lemma}[Axis and angle of an $\mathrm{SO}(3)$ element]
    \label{lem:so3-axis-angle}
    \uses{def:pure-euclidean}
    Every $A \in \mathrm{SO}(3)$ admits a unit axis $u$ with $Au = u$ and an angle
    $\varphi$ such that, on the orthogonal complement $u^\perp$ with respect to an
    orthonormal frame $\{u, w, u\times w\}$, $A$ acts as the planar rotation by
    $\varphi$. Writing $\varphi = 2\theta$, the data is a unit axis $u$ together
    with $\theta$.
\end{lemma}
\begin{proof}
    \uses{def:pure-euclidean}
    By \texttt{GroupTheory.SO3.exists\_stationary\_vec} there is a unit vector $u$
    with $Au = u$. Since $A$ is orthogonal it preserves $u^\perp$, and its
    restriction to the $2$-dimensional space $u^\perp$ is an element of
    $\mathrm{SO}(2)$, hence a planar rotation by a unique angle $\varphi$; set
    $\theta = \varphi/2$.
\end{proof}

\begin{lemma}[Conjugation fixes its axis]
    \label{lem:rotation-fixes-axis}
    \uses{lem:conj-linear, lem:pure-square}
    For a unit pure quaternion $u$ and $\theta \in \mathbb{R}$, set
    $\alpha = \cos\theta + (\sin\theta)\,u \in \mathbb{H}^1$ (using $u^2 = -1$,
    \ref{lem:pure-square}). Then $\rho_\alpha(u) = u$.
\end{lemma}
\begin{proof}
    \uses{lem:pure-square}
    Since $u^2 = -1$ (\ref{lem:pure-square}), $u$ commutes with
    $\alpha = \cos\theta + (\sin\theta) u$ (a polynomial in $u$), so
    $\alpha u = u\alpha$ and hence $\rho_\alpha(u) = \alpha u\alpha^{-1} = u$.
\end{proof}

\begin{lemma}[Conjugation rotates the orthogonal plane by $2\theta$]
    \label{lem:rotation-on-perp}
    \uses{lem:conj-linear, lem:pure-square, lem:pure-orthogonal, lem:pure-product}
    With $u, \alpha$ as in \ref{lem:rotation-fixes-axis}, for any unit pure $w$
    orthogonal to $u$ the map $\rho_\alpha$ acts on the plane
    $\operatorname{span}\{w, u\times w\}$ as the rotation by angle $2\theta$:
    $\rho_\alpha(w) = (\cos 2\theta)\,w + (\sin 2\theta)\,(u\times w)$.
\end{lemma}
\begin{proof}
    \uses{lem:pure-square, lem:pure-orthogonal, lem:pure-product}
    Since $u\perp w$ we have $uw = u\times w$ and $wu = -uw$
    (\ref{lem:pure-orthogonal}, \ref{lem:pure-product}). Expanding
    $\rho_\alpha(w) = \alpha w\bar\alpha$ with $\alpha = \cos\theta + (\sin\theta)u$
    and using $u^2 = -1$ (\ref{lem:pure-square}) gives, by the double-angle
    formulas, $(\cos^2\theta - \sin^2\theta)w + 2\sin\theta\cos\theta\,(u\times w)
    = (\cos 2\theta)w + (\sin 2\theta)(u\times w)$.
\end{proof}

\begin{lemma}[Rotation map is surjective]
    \label{lem:rotation-surjective}
    \uses{prop:rotation, lem:so3-axis-angle, lem:rotation-fixes-axis, lem:rotation-on-perp, lem:pure-square, lem:pure-orthogonal}
    Every $A \in \mathrm{SO}(3)$ equals $\rho_\alpha$ for some
    $\alpha \in \mathbb{H}^1$.
\end{lemma}
\begin{proof}
    \uses{lem:so3-axis-angle, lem:rotation-fixes-axis, lem:rotation-on-perp, lem:pure-square, lem:pure-orthogonal}
    By \ref{lem:so3-axis-angle}, $A$ has a unit axis $u$ and acts on $u^\perp$ as
    rotation by $2\theta$. Set $\alpha = \cos\theta + (\sin\theta)u \in \mathbb{H}^1$.
    Then $\rho_\alpha(u) = u = Au$ (\ref{lem:rotation-fixes-axis}) and on $u^\perp$
    both $\rho_\alpha$ and $A$ are the rotation by $2\theta$
    (\ref{lem:rotation-on-perp}). As $\rho_\alpha$ and $A$ agree on the basis
    $\{u, w, u\times w\}$, they are equal.
\end{proof}

\begin{corollary}[The double cover $\mathbb{H}^1 \to \mathrm{SO}(3)$]
    \label{cor:cover}
    \uses{prop:rotation, lem:rotation-kernel, lem:rotation-surjective}
    The map $\alpha \mapsto \rho_\alpha$ is a surjective group homomorphism
    $\mathbb{H}^1 \to \mathrm{SO}(3)$ with kernel $\{\pm 1\}$, giving the exact
    sequence
    \[
      1 \to \{\pm 1\} \to \mathbb{H}^1 \to \mathrm{SO}(3) \to 1.
    \]
\end{corollary}
\begin{proof}
    \uses{prop:rotation, lem:rotation-kernel, lem:rotation-surjective}
    The assignment is a homomorphism since
    $\rho_{\alpha\beta}(v) = (\alpha\beta) v (\alpha\beta)^{-1}
    = \alpha(\beta v \beta^{-1})\alpha^{-1} = (\rho_\alpha \circ \rho_\beta)(v)$,
    and lands in $\mathrm{SO}(3)$ by \ref{prop:rotation}. Surjectivity is
    \ref{lem:rotation-surjective} and the kernel is $\{\pm 1\}$ by
    \ref{lem:rotation-kernel}.
\end{proof}

\end{document}
